\section*{Erd\H{o}s Problem 1209 --- GPT 6 Astra Ultra partial attempt}


\subsection*{FORMAL RESTATEMENT}
For a strictly increasing integer sequence $(a_k)$ with sufficiently rapid growth, suppose one translate $n+(a_k)$ consists entirely of primes. Must infinitely many translates do so? Ask the same question with squarefree values in place of primes. Separately, for the sequence $2^{2^k}$, does some integer $n$ make $n+2^{2^k}$ prime for every $k$, squarefree for every $k$, prime for infinitely many $k$, or squarefree for infinitely many $k$?

\subsection*{CONTEXT AND SCOPE}
This replaces the earlier statement-only record with an elementary mathematical attempt. The arguments below establish only their stated partial results; no novelty or full solution is claimed.

The source does not quantify sufficiently rapid growth or specify the allowed shifts n and initial k. Its remarks give counterexamples to both general translate implications and discuss a negative all-prime special case; the remaining subquestions explain its open status. Preserve these as distinct questions rather than claiming they all remain unresolved.

\subsection*{ATTACK PLAN AND WORK}
There are rigorous congruence restrictions for the special sequence. For every
even shift $n$, the values $n+2^{2^k}$ are even and greater than two for large
$k$, so at most finitely many are prime. For $k\ge1$,
$2^{2^k}\equiv1\pmod3$; hence $n\equiv2\pmod3$ also rules out infinitely many
prime values. For $k\ge2$, the analogous congruence modulo five excludes
$n\equiv4\pmod5$. In each case the finitely many possible exceptions where
the value equals the dividing prime are kept separate.

For the shift $n=1$, put $F_k=2^{2^k}+1$. The identity
\[\prod_{j=0}^{k-1}F_j=F_k-2\tag{1}\]
follows inductively from $(u-1)(u+1)=u^2-1$.
It implies pairwise coprimality of the $F_k$: a common divisor of $F_j$ and
$F_k$ for $j<k$ divides two, while both numbers are odd. Thus infinitely many
different prime divisors occur among these terms.

That conclusion says nothing about infinitely many prime terms or squarefree
terms. Pairwise coprime numbers can all be composite, and can all have square
factors. It would be invalid to promote (1) to either requested property.
The source's already negative general translate implications are likewise
not being reproved or contradicted by these special-sequence observations.

\subsection*{VERIFICATION AND REMAINING GAP}
Resolve the remaining prime and squarefree translate questions with the
intended range of shifts and growth hypotheses specified. Congruence exclusions
and pairwise coprimality leave the essential existence questions unsettled.

\subsection*{SOURCES}
https://www.erdosproblems.com/1209

\subsection*{FINAL}
\textbf{UNRESOLVED.} The full source problem remains outside the scope of the proved partial results.

\subsection*{COMPLETION ESTIMATE}
$6\%$. This is a conservative subjective estimate toward the whole problem, not a probability that the argument is correct or a claim that a fixed fraction of the problem has been solved.
